% GENERATED FILE. Edit canonical lessons, then run npm run generate:books.
% canonical-source-hash: fnv1a64:a372c21c21eca291
% canonical-lessons: TE-C110-talam, TE-C110-talamcevi, TE-C110-tadu, TE-C110-meku, TE-C110-kattera

\chapter{Lock, Key, Rope, Nail, Scissors}
\label{ch:te-lock-key-rope-nail-scissors}
\hlchaptermodality{110}

\begin{chapteropening}
\textbf{By the end of this chapter:} \emph{I can say a lock, a key, a rope, a nail and scissors.}

Complete the last lesson of chapter 110: I can say a lock, a key, a rope, a nail and scissors.
\end{chapteropening}

\section[\texorpdfstring{\te{తాళం} (tāḷaṁ)}{తాళం (tāḷaṁ)}]{\te{తాళం} (tāḷaṁ) — a lock}

\label{lesson:TE-C110-talam}



\begin{quote}
Before the new one: say the Telugu for a blanket, then the Telugu for a curtain, a screen.
\end{quote}

\subsection*{You'll want to know: \te{తాళం}}
\textbf{\te{తాళం}} — \emph{tāḷaṁ} — \textquotedblleft{}a lock\textquotedblright{}.

\begin{grammarlens}[title={What you've built}]
One more word for this chapter.
\end{grammarlens}

\subsection*{Your turn}
\begin{itemize}
\raggedright
  \item \emph{Say it:} \emph{tāḷaṁ}
  \item \emph{Say it:} \emph{tāḷaṁ}, once more
  \item \emph{Say it:} say \emph{tāḷaṁ}
  \item \emph{Recall:} say the Telugu for a blanket, then the Telugu for a curtain, a screen, then say \emph{tāḷaṁ} again
\end{itemize}

\begin{tcolorbox}[breakable,colback=teal!4,colframe=teal!35!black,title={Before you move on}]
What does \te{తాళం} mean? (\textquotedblleft{}A lock\textquotedblright{}.) Say it once more.
\end{tcolorbox}
\section[\texorpdfstring{\te{తాళంచెవి} (tāḷaṁcevi)}{తాళంచెవి (tāḷaṁcevi)}]{\te{తాళంచెవి} (tāḷaṁcevi) — a key}

\label{lesson:TE-C110-talamcevi}



\begin{quote}
Before the new one: say the Telugu for a curtain, a screen, then the Telugu for a lock.
\end{quote}

\subsection*{You'll want to know: \te{తాళంచెవి}}
\textbf{\te{తాళంచెవి}} — \emph{tāḷaṁcevi} — \textquotedblleft{}a key\textquotedblright{}.

\begin{grammarlens}[title={What you've built}]
One more word for this chapter.
\end{grammarlens}

\subsection*{Your turn}
\begin{itemize}
\raggedright
  \item \emph{Say it:} \emph{tāḷaṁcevi}
  \item \emph{Say it:} \emph{tāḷaṁcevi}, once more
  \item \emph{Say it:} say \emph{tāḷaṁcevi}
  \item \emph{Recall:} say the Telugu for a curtain, a screen, then the Telugu for a lock, then say \emph{tāḷaṁcevi} again
\end{itemize}

\begin{tcolorbox}[breakable,colback=teal!4,colframe=teal!35!black,title={Before you move on}]
What does \te{తాళంచెవి} mean? (\textquotedblleft{}A key\textquotedblright{}.) Say it once more.
\end{tcolorbox}
\section[\texorpdfstring{\te{తాడు} (tāḍu)}{తాడు (tāḍu)}]{\te{తాడు} (tāḍu) — a rope}

\label{lesson:TE-C110-tadu}



\begin{quote}
Before the new one: say the Telugu for a lock, then the Telugu for a key.
\end{quote}

\subsection*{You'll want to know: \te{తాడు}}
\textbf{\te{తాడు}} — \emph{tāḍu} — \textquotedblleft{}a rope\textquotedblright{}.

\begin{grammarlens}[title={What you've built}]
One more word for this chapter.
\end{grammarlens}

\subsection*{Your turn}
\begin{itemize}
\raggedright
  \item \emph{Say it:} \emph{tāḍu}
  \item \emph{Say it:} \emph{tāḍu}, once more
  \item \emph{Say it:} say \emph{tāḍu}
  \item \emph{Recall:} say the Telugu for a lock, then the Telugu for a key, then say \emph{tāḍu} again
\end{itemize}

\begin{tcolorbox}[breakable,colback=teal!4,colframe=teal!35!black,title={Before you move on}]
What does \te{తాడు} mean? (\textquotedblleft{}A rope\textquotedblright{}.) Say it once more.
\end{tcolorbox}
\section[\texorpdfstring{\te{మేకు} (mēku)}{మేకు (mēku)}]{\te{మేకు} (mēku) — a nail (to hammer)}

\label{lesson:TE-C110-meku}



\begin{quote}
Before the new one: say the Telugu for a key, then the Telugu for a rope.
\end{quote}

\subsection*{You'll want to know: \te{మేకు}}
\textbf{\te{మేకు}} — \emph{mēku} — \textquotedblleft{}a nail (to hammer)\textquotedblright{}.

\begin{grammarlens}[title={What you've built}]
One more word for this chapter.
\end{grammarlens}

\subsection*{Your turn}
\begin{itemize}
\raggedright
  \item \emph{Say it:} \emph{mēku}
  \item \emph{Say it:} \emph{mēku}, once more
  \item \emph{Say it:} say \emph{mēku}
  \item \emph{Recall:} say the Telugu for a key, then the Telugu for a rope, then say \emph{mēku} again
\end{itemize}

\begin{tcolorbox}[breakable,colback=teal!4,colframe=teal!35!black,title={Before you move on}]
What does \te{మేకు} mean? (\textquotedblleft{}A nail (to hammer)\textquotedblright{}.) Say it once more.
\end{tcolorbox}
\section[\texorpdfstring{\te{కత్తెర} (katterā)}{కత్తెర (katterā)}]{\te{కత్తెర} (katterā) — scissors}

\label{lesson:TE-C110-kattera}



\begin{quote}
Before the new one: say the Telugu for a rope, then the Telugu for a nail.
\end{quote}

\subsection*{You'll want to know: \te{కత్తెర}}
\textbf{\te{కత్తెర}} — \emph{katterā} — \textquotedblleft{}scissors\textquotedblright{}.

\begin{grammarlens}[title={What you've built}]
One more word for this chapter.
\end{grammarlens}

\subsection*{Your turn}
\begin{itemize}
\raggedright
  \item \emph{Say it:} \emph{katterā}
  \item \emph{Say it:} \emph{katterā}, once more
  \item \emph{Say it:} say \emph{katterā}
  \item \emph{Recall:} say the Telugu for a rope, then the Telugu for a nail, then say \emph{katterā} again
\end{itemize}

\begin{tcolorbox}[breakable,colback=teal!4,colframe=teal!35!black,title={Before you move on}]
What does \te{కత్తెర} mean? (\textquotedblleft{}Scissors\textquotedblright{}.) Say it once more.
\end{tcolorbox}
